\section{The binomial packet and phase-uniform anti-cancellation}
\label{sec:terminal-modal-packet}

Define the ideal entry packet
\begin{equation}
 g_j=-\frac{\alpha}{2}(1-q)^m b_{m,j},
 \qquad b_{m,j}=2^{-m}\binom mj.
 \label{eq:terminal-modal-packet}
\end{equation}
Its center height is
\begin{equation}
 \frac{\max_j|g_j|}{q^{5/2}}
 \longrightarrow 2e^{-1/16}\sqrt{\frac2\pi}.
 \label{eq:terminal-modal-packet-height}
\end{equation}

\begin{lemma}[Exact packet coefficients; proved here]
\label{lem:terminal-modal-packet-transform}
For $1\le h\le m-1$, the weighted cosine coefficient of $g$ is
\begin{equation}
 \widehat g_h=-\frac{\alpha(1-q)^m}{m}
 \left[
  \cos^m\phi_h\cos\frac{h\pi}{2}
  -2^{-m-1}(1+(-1)^h)
 \right].
 \label{eq:terminal-modal-packet-transform}
\end{equation}
In particular, for $h=2s$,
\begin{equation}
 G_{2s}:=\widehat g_{2s}
 =-\frac{\alpha(1-q)^m}{m}
 \left[(-1)^s\cos^m\frac{s\pi}{m}-2^{-m}\right].
 \label{eq:terminal-modal-even-packet}
\end{equation}
For $s=o(\sqrt m)$,
$|G_{2s}|\sim e^{-1/16}\alpha/m=16e^{-1/16}q^3$.
\end{lemma}

\begin{proof}
If $J\sim\operatorname{Binomial}(m,1/2)$, then
\[
 \mathbb E e^{itJ}=e^{imt/2}\cos^m(t/2),
 \qquad
 \mathbb E\cos(tJ)=\cos(mt/2)\cos^m(t/2).
\]
The degree weights double every interior binomial mass and leave the two
endpoint masses single.  Taking $t=h\pi/m$, subtracting those endpoint
copies, multiplying by $-\alpha(1-q)^m/2$, and dividing by $N_h=m$ gives
\eqref{eq:terminal-modal-packet-transform}.  The even-mode specialization and
asymptotic follow immediately.
\end{proof}

\begin{lemma}[Exact position/velocity propagators; proved here]
\label{lem:terminal-modal-propagators}
Put $L=(I+P_{\rm row})/2$ and
$s_k=(1-q)^{-k}r_k^x$.  Every projection-inactive full-face trajectory obeys
\begin{equation}
 s_{k+1}=2Ls_k-Ls_{k-1}.
 \label{eq:terminal-modal-scaled-recurrence}
\end{equation}
Let $K_0=I,K_1=L$ and $J_0=0,J_1=I$ be the two fundamental solutions of
this recurrence.  On the mode with $\theta_h=h\pi/m$, their multipliers are
\begin{align}
 K_k(\theta_h)&=\cos^k\frac{\theta_h}{2}
                         \cos\frac{k\theta_h}{2},
 \label{eq:terminal-modal-position-kernel}\\
 J_k(\theta_h)&=\cos^{k-1}\frac{\theta_h}{2}
 \frac{\sin(k\theta_h/2)}{\sin(\theta_h/2)},
 \qquad k\ge1,
 \label{eq:terminal-modal-velocity-kernel}
\end{align}
with the second expression continued at $\theta_h=0$.  Both kernels are
entrywise nonnegative on the endpoint path.  The row sums of $K_k$ and $J_k$
are respectively $1$ and $k$.

More precisely, if $\mathsf T_j$ is the Chebyshev polynomial, then
\begin{align}
 K_k&=2^{-k}\sum_{j=0}^k\binom{k}{j}\mathsf T_j(P_{\rm row}),
 \label{eq:terminal-modal-position-expansion}\\
 J_k&=2^{-(k-1)}\sum_{a,b=0}^{k-1}\binom{k-1}{b}
 \mathsf T_{|k-1-a-b|}(P_{\rm row}).
 \label{eq:terminal-modal-velocity-expansion}
\end{align}
Consequently, with
\begin{equation}
 c=r_0^x-g,\qquad
 w=\frac{r_1^x}{1-q}-Lr_0^x,
 \label{eq:terminal-modal-propagator-data}
\end{equation}
one has the exact decomposition
\begin{equation}
 r_k^x=(1-q)^k\left(K_kg+K_kc+J_kw\right).
 \label{eq:terminal-modal-propagator-decomposition}
\end{equation}
For $1\le h\le m-1$, the mode-$h$ coefficient of $w$ is
$\cos\phi_h\sin\phi_hD_h=q\cos^2\phi_hV_h$.
\end{lemma}

\begin{proof}
Since $B=(1-q^2)L$, rescaling
\eqref{eq:terminal-modal-second-order} gives
\eqref{eq:terminal-modal-scaled-recurrence}.  Solving its scalar recurrence
at $L$-eigenvalue $\cos^2(\theta_h/2)$ gives
\cref{eq:terminal-modal-position-kernel,eq:terminal-modal-velocity-kernel}.
The position expansion follows by taking the real part of
$2^{-k}(1+e^{i\theta})^k$.  For the velocity expansion, multiply
\[
 \frac{\sin(kx)}{\sin x}=\sum_{a=0}^{k-1}e^{i(k-1-2a)x}
\]
by the binomial expansion of $\cos^{k-1}x$, set $x=\theta/2$, and take real
parts.

To interpret the operator expansions, extend a path vector evenly to the cycle
$\mathbb Z/(2m)$.  On the endpoint-weighted cosine basis, multiplication by
$\cos(j\theta_h)$, equivalently $\mathsf T_j(P_{\rm row})$, is exactly the
average of the two folded shifts by $\pm j$.  It is positivity preserving and
has row sum one, including all endpoint reflections and aliases.  The
coefficient sums in \eqref{eq:terminal-modal-position-expansion} and
\eqref{eq:terminal-modal-velocity-expansion} are $1$ and $k$.  This proves the
kernel claims.  The fundamental-solution initial conditions give
\eqref{eq:terminal-modal-propagator-decomposition}.  The formula for $w_h$
follows from \eqref{eq:terminal-modal-quadratures}, and its second form follows
from \cref{lem:terminal-modal-entry-velocity}.
\end{proof}

\begin{lemma}[Ideal packet remains nonpositive; proved here]
\label{lem:terminal-modal-ideal-sign}
The ideal residual trajectory $r_k^{\rm id}=(1-q)^kK_kg$ is coordinatewise
nonpositive at every integer time $k\ge0$.
\end{lemma}

\begin{proof}
The packet $g$ is coordinatewise nonpositive, while
\cref{lem:terminal-modal-propagators} shows that $K_k$ is entrywise
nonnegative.
\end{proof}

In particular, the ideal residual has zero safe-envelope correction at every
time.  Moreover $K_k$ is an $\ell_\infty$ contraction, but the direct bound on
$J_kw$ loses the factor $k$ because $J_k$ has row sum $k$.  Thus the precise
remaining obstruction on the regime side is signed or variation cancellation
in the velocity remainder $w$; the sign fact alone does not prove raw-point
positivity or close \cref{lem:terminal-modal-regime}.

\begin{lemma}[Uniform cosine-square floor; proved here]
\label{lem:terminal-modal-cosine-square}
For every integer $H\ge8$ and every real $x$,
\begin{equation}
 \sum_{s=1}^{H}\cos^2(sx)\ge\frac{H}{16}.
 \label{eq:terminal-modal-cosine-square}
\end{equation}
\end{lemma}

\begin{proof}
If $\sin x=0$, the sum equals $H$.  Otherwise, the finite sum identity gives
\[
 \sum_{s=1}^H\cos^2(sx)
 =\frac H2+\frac{\sin(Hx)\cos((H+1)x)}{2\sin x}.
\]
If $|\sin x|\ge2/H$, the right side is at least $H/4$.  Otherwise let
$\xi$ be the distance from $x$ to $\pi\mathbb Z$.  Since
$\sin(\pi/H)\ge2/H$, one has $\xi<\pi/H$.  For
$1\le s\le\lfloor H/4\rfloor$, $s\xi\le\pi/4$, so each corresponding
square is at least $1/2$.  Since
$\lfloor H/4\rfloor\ge H/8$ for $H\ge8$, the result follows.
\end{proof}

The position half of the former entry-profile interface is now
\cref{lem:terminal-modal-position-profile}.  The velocity half remains open.
It is displayed as a lemma to freeze the needed inequality, but is not
asserted as proved.

\begin{lemma}[Entry velocity profile; open]
\label{lem:terminal-modal-packet}
For all sufficiently large $m$, the literal full-face entry coefficients
satisfy, for $1\le s\le H_m$,
\begin{equation}
 \frac m\alpha\left|V_{2s}+\frac{G_{2s}}5\right|\le\frac14.
 \label{eq:terminal-modal-entry-velocity-profile}
\end{equation}
\end{lemma}

The second profile is deliberately centered at $-G_{2s}/5$, not at zero.
The screen below indicates a nonvanishing smooth correction, so an $o(1)$
profile error is neither needed nor suggested.

\begin{lemma}[Profiles imply the quadrature target; proved here]
\label{lem:terminal-modal-profile-implies-quadrature}
For all sufficiently large $m$, the bounds in
\eqref{eq:terminal-modal-entry-position-profile} and
\eqref{eq:terminal-modal-entry-velocity-profile}, together with inactivity
of the first full-face projection, imply
\begin{equation}
 |C_{2s}-G_{2s}|+|D_{2s}|\le\frac{|G_{2s}|}{64},
 \qquad 1\le s\le H_m.
 \label{eq:terminal-modal-entry-target}
\end{equation}
\end{lemma}

\begin{proof}
Put $g_s=m|G_{2s}|/\alpha$.  Uniformly on the stated band,
$\cos^m(s\pi/m)-2^{-m}\ge14/15$ for all sufficiently large $m$.
Also $(1-q)^m\ge1-mq=15/16$.  Thus $g_s\ge7/8$.
By \cref{lem:terminal-modal-entry-velocity} and $\cot x\le1/x$,
\[
 \frac{|C_{2s}-G_{2s}|+|D_{2s}|}{|G_{2s}|}
 \le \frac1{256g_s}
 +q\cot\frac{s\pi}{m}\left(\frac15+\frac1{4g_s}\right)
 \le\frac1{224}+\frac{17}{560\pi}<\frac1{64}.
\]
This is exactly \eqref{eq:terminal-modal-entry-target}.
\end{proof}

\begin{lemma}[Projection and global-envelope regime; open]
\label{lem:terminal-modal-regime}
Put
\begin{equation}
 K_m=\left\lfloor\frac1{8q}\log\frac1q\right\rfloor.
 \label{eq:terminal-modal-horizon}
\end{equation}
For every $1\le k\le K_m$, the raw point
$y_{k-1}-r(y_{k-1})$ is strictly positive, and the resulting full-face state
satisfies
\begin{equation}
 \min_j p_{k,j}>\delta_k.
 \label{eq:terminal-modal-unclipped-target}
\end{equation}
\end{lemma}

This statement preserves both nonlinearities that the full-face linear
recurrence omits: the coordinate projection and the moving global correction.
The separate proper-prefix chronology lemma preserves the same issues during
support growth.

\begin{theorem}[Conditional terminal logarithmic block; proved here]
\label{thm:terminal-modal-conditional-log}
Assume the conclusions of
\cref{lem:terminal-modal-packet,lem:terminal-modal-regime}.
Then, for
all sufficiently large $m$, the named execution has no terminal one-sided
certificate during its first $K_m$ full-face steps.  Hence its terminal block
has length
\begin{equation}
 \Omega\!\left(q^{-1}\log\frac1q\right).
 \label{eq:terminal-modal-conditional-log}
\end{equation}
This conclusion is for the named literal recurrence only.
\end{theorem}

\begin{proof}
By
\cref{lem:terminal-modal-position-profile,lem:terminal-modal-profile-implies-quadrature},
the proved position profile and the open velocity profile give
\eqref{eq:terminal-modal-entry-target}.
For $1\le s\le H_m$, let $A_s=|G_{2s}|$.  Uniformly on this band and for all
sufficiently large $m$,
\begin{equation}
 \frac{\alpha}{3m}\le A_s\le\frac{2\alpha}{m}.
 \label{eq:terminal-modal-amplitude-bounds}
\end{equation}
Indeed, $(1-q)^m\ge1-mq=15/16$, while
$m(s\pi/m)^2\le\pi^2/(64\log(16m))$ and
$\log\cos z\ge-z^2$ for sufficiently small $z$.

Fix $1\le k\le K_m$.  For $h=2s$, the phase in
\cref{eq:terminal-modal-solution} is $\phi_{2s}=s\pi/m$.  Moreover,
\begin{equation}
 \cos^{k}\frac{H_m\pi}{m}
 \ge \exp\!\left(-\frac{\pi^2}{32}\right)>\frac12.
 \label{eq:terminal-modal-spatial-damping}
\end{equation}
Apply \cref{lem:terminal-modal-cosine-square} at
$x=k\pi/m$.  The ideal part of the band coefficient vector has Euclidean norm
at least
\[
 (1-q)^k\cdot\frac{\alpha}{3m}\cdot\frac12
 \sqrt{\frac{H_m}{16}}
 =\frac{\alpha(1-q)^k\sqrt{H_m}}{24m}.
\]
By \eqref{eq:terminal-modal-entry-target}, its perturbation has norm at most
\[
 (1-q)^k\cdot\frac{2\alpha}{m}\cdot
 \frac{\sqrt{H_m}}{64}
 =\frac{\alpha(1-q)^k\sqrt{H_m}}{32m}.
\]
The reverse triangle inequality therefore gives
\begin{equation}
 \left(\sum_{s=1}^{H_m}a_{2s,k}^2\right)^{1/2}
 \ge\frac{\alpha(1-q)^k\sqrt{H_m}}{96m}.
 \label{eq:terminal-modal-band-energy}
\end{equation}
Parseval and \cref{eq:terminal-modal-popoviciu} imply
\begin{equation}
 R_k\ge
 \frac{\sqrt2\alpha(1-q)^k\sqrt{H_m}}{96m}
 =\frac{\sqrt2}{6}q^3(1-q)^k\sqrt{H_m}.
 \label{eq:terminal-modal-range-lower}
\end{equation}

Because $q\le1/32$,
\[
 (1-q)^{K_m}
 \ge \exp\!\left(-\frac{qK_m}{1-q}\right)
 \ge q^{4/31}.
\]
Also
$\sqrt{H_m}=\Omega(q^{-1/4}(\log(1/q))^{-1/4})$.  Consequently the ratio of
the right side of \eqref{eq:terminal-modal-range-lower} to $q^3$ diverges
like a positive constant times
$q^{-15/124}(\log(1/q))^{-1/4}$.  It is therefore greater than $1/5$ for all
sufficiently small $q$.  The open regime lemma permits
\cref{lem:terminal-modal-range}, which proves non-certification through
$K_m$.
\end{proof}

\begin{remark}[Why one mode is insufficient]
For each fixed $h$, the ideal coefficient is only $\Theta(q^3)$, exactly the
certificate scale, and it can pass through zero.  The gain comes from a band
of $H_m$ orthogonal modes and the phase-uniform
\cref{lem:terminal-modal-cosine-square}; it is not a surviving-mode argument.
The conservative proof constant $1/8$ is not intended to predict the measured
first-certificate constant, which is closer to
$\frac12\log(1/q)$ over the screened sizes.
\end{remark}
